\section{Valid Parentheses}
\label{sec:sq-valid-parentheses}

\problemstatement{Given a string $s$ consisting only of the characters
\texttt{'('}, \texttt{')'}, \texttt{'['}, \texttt{']'}, \texttt{'\{'},
\texttt{'\}'}, determine whether the brackets are \textbf{valid}: every
closing bracket must match the most recently opened, still-unmatched
bracket of the same type, and every opening bracket must eventually be
closed.}

\begin{patternbox}
\emph{``Matching/nesting''} is the most direct stack keyword there is: the
defining feature of valid nesting is that the bracket needing to be closed
\textbf{next} is always the \textbf{most recently opened} one still
unmatched -- exactly LIFO order. Whenever a problem describes validity in
terms of nested or paired structures (brackets, tags, nested function
calls), a stack of ``currently open, unmatched'' items is almost always
the right tool.
\end{patternbox}

\subsection*{Understanding the problem carefully}

Scan the string left to right. Every opening bracket is a commitment that
must be fulfilled later, in reverse order of when the commitments were
made -- the bracket opened most recently must be the first one closed,
because brackets cannot cross (that would not be ``nesting''). Pushing
each opening bracket onto a stack, and popping when we see a closing
bracket, mirrors this exactly: the stack's top is always the innermost
currently-open bracket, which is precisely the one any new closing
bracket must match.

\begin{ideabox}[Stack of Open Brackets]
For each character: if it is an opening bracket, push it. If it is a
closing bracket: the stack must be non-empty, and its top must be the
\emph{matching} opening bracket; if either check fails, the string is
invalid. Otherwise, pop the stack (this pair is now resolved). After
scanning the whole string, the string is valid if and only if the stack
is empty (every opened bracket was eventually closed).
\end{ideabox}

\subsection*{Why both checks (stack non-empty, and top matches) are necessary}

If the stack is empty when a closing bracket appears, there is no
open bracket left to match it against -- the closing bracket is
unmatched, and the string is invalid (e.g.\ \texttt{")"} alone). If the
stack's top does not match the closing bracket's type (e.g.\ the top is
\texttt{'('} but we just saw \texttt{']'}), the brackets have crossed
incorrectly (e.g.\ \texttt{"(]"}), which also invalidates the string. And
even if every individual closing bracket finds a match, any \emph{leftover}
opening brackets at the end (a non-empty stack) mean some opens were never
closed at all (e.g.\ \texttt{"(("}, missing both closing parentheses) --
this is why the final emptiness check is also required, not just the
per-character checks.

\subsection*{Algorithm}

\begin{enumerate}
  \item Initialize an empty stack.
  \item For each character $c$ in $s$:
  \begin{enumerate}
    \item If $c$ is an opening bracket: push $c$.
    \item Else ($c$ is a closing bracket): if the stack is empty, or its
          top is not the matching opening bracket for $c$, return
          \texttt{false}. Otherwise, pop the stack.
  \end{enumerate}
  \item Return whether the stack is now empty.
\end{enumerate}

\begin{algorithm}[H]
\caption{Valid Parentheses}
\begin{algorithmic}[1]
\Procedure{IsValid}{$s$}
  \State $\textit{stack} \gets$ empty
  \For{each character $c$ in $s$}
    \If{$c \in \{(, [, \{\}$}
      \State push $c$ onto \textit{stack}
    \Else
      \If{\textit{stack} is empty \textbf{or} top of \textit{stack} does not match $c$}
        \State \Return \textbf{false}
      \EndIf
      \State pop \textit{stack}
    \EndIf
  \EndFor
  \State \Return (\textit{stack} is empty)
\EndProcedure
\end{algorithmic}
\end{algorithm}

\subsection*{Dry run}

\dryrun{Take $s = \texttt{"\{[()]\}"}$.}

\begin{itemize}
  \item \texttt{'\{'}: push. Stack $=[\{]$.
  \item \texttt{'['}: push. Stack $=[\{,[]$.
  \item \texttt{'('}: push. Stack $=[\{,[,(]$.
  \item \texttt{')'}: top is \texttt{'('}, matches. Pop. Stack
        $=[\{,[]$.
  \item \texttt{']'}: top is \texttt{'['}, matches. Pop. Stack $=[\{]$.
  \item \texttt{'\}'}: top is \texttt{'\{'}, matches. Pop. Stack $=[]$.
\end{itemize}

Stack is empty at the end: \texttt{true}.

Now take $s = \texttt{"(]"}$: push \texttt{'('}; then see \texttt{']'},
top is \texttt{'('}, which does not match \texttt{']'} -- return
\texttt{false} immediately.

\subsection*{C++ implementation}

\begin{lstlisting}
#include <bits/stdc++.h>
using namespace std;

bool isValid(string s) {
    stack<char> st;
    for (char c : s) {
        if (c == '(' || c == '[' || c == '{') {
            st.push(c);
        } else {
            if (st.empty()) return false;
            char top = st.top();
            if ((c == ')' && top != '(') ||
                (c == ']' && top != '[') ||
                (c == '}' && top != '{')) {
                return false;
            }
            st.pop();
        }
    }
    return st.empty();
}

int main() {
    cout << (isValid("{[()]}") ? "true" : "false") << endl;
    cout << (isValid("(]") ? "true" : "false") << endl;
    return 0;
}
\end{lstlisting}

\begin{complexitybox}
\textbf{Time Complexity.} Each character is processed once with constant
work, giving
\[
  O(n).
\]
\textbf{Space Complexity.} $O(n)$ in the worst case (a string of all
opening brackets pushes every character onto the stack).
\end{complexitybox}

\begin{takeawaybox}
``The next thing that needs to be resolved is always the most recently
opened, still-unresolved thing'' is the defining signature of a stack
problem. This is not limited to brackets -- the same template applies to
validating nested tags, matching function call/return pairs, or undoing
actions in reverse order of when they were performed.
\end{takeawaybox}
